\section{Results}

This section summarizes the main results from the thesis in the same order as
the experimental workflow. The first result group reports the xLSTM
model-selection stage, where validation loss was used to choose a recurrent
checkpoint and to estimate the practical compute budget of the search. The
second result group reports the matched xLSTM, GPT-2-style Transformer, and
LLaMA-style Transformer comparison on next-token metrics. The third group
summarizes the generated-MIDI evaluation, and the final group reports the
interactive Gradio deliverable.

\subsection{xLSTM Model-Selection Results}

The xLSTM model-selection results are validation-stage results. They show that
the implemented training pipeline can train xLSTM language models on REMI-style
music tokens, compare checkpoints, and identify stronger configurations for the
controlled Transformer comparison. They should therefore be interpreted as
selection evidence, not as a final statement about musical quality.

\subsubsection{Validation Results}

Table~\ref{tab:result-checkpoints} summarizes representative completed xLSTM
runs from the local \verb|summary.json| files. The strongest validation-loss
checkpoint is the base xLSTM with tokenizer 14 and batch size 16. The
batch-size-8 tokenizer-14 run is almost tied and keeps the same practical setup
used by the broader sweep. The medium model does not improve the
validation-selected result under its current 15-epoch restart schedule, even
though it has more than twice as many parameters.

\begin{table}[htbp]
	\centering
	\scriptsize
	\resizebox{\textwidth}{!}{%
		\begin{tabular}{lrrrrrrr}
			\toprule
			Run                                & Params & Tok. & Batch & Steps & Best step & Best val. loss  & Best val. top-5 acc. \\
			\midrule
			\texttt{tok\_14\_batch\_16}        & 3.70M  & 14   & 16    & 7287  & 4000      & \textbf{1.4456} & \textbf{0.8619}      \\
			\texttt{tok\_14}                   & 3.70M  & 14   & 8     & 14574 & 6000      & 1.4471          & 0.8614               \\
			\texttt{constant\_1e-4}            & 3.67M  & 11   & 8     & 12705 & 12000     & 1.5759          & 0.8365               \\
			\texttt{xlstm-small-batch8-tok11}  & 1.63M  & 11   & 8     & 12705 & 12000     & 1.5698          & 0.8388               \\
			\texttt{xlstm-base-batch8-tok11}   & 3.67M  & 11   & 8     & 12705 & 10000     & 1.5811          & 0.8367               \\
			\texttt{xlstm-medium-batch8-tok14} & 9.28M  & 14   & 8     & 72870 & 9000      & 1.5173          & 0.8542               \\
			\bottomrule
		\end{tabular}
	}
	\caption{Representative xLSTM validation results from local experiment summaries. Lower validation loss is better. The medium row reports the early best checkpoint, not the final checkpoint.}
	\label{tab:result-checkpoints}
\end{table}

The validation ranking supports two main observations. First, tokenizer 14 is a
better downstream representation for the current xLSTM setup than tokenizer 11,
even though tokenizer 11 was the best tokenizer-only round-trip candidate. This
confirms that compactness and reconstruction metrics are useful diagnostics but
do not fully predict language-model training behavior. Second, simply increasing
model size is not enough. The medium tokenizer-14 run reaches its best
validation loss early and then overfits, so its final checkpoint is worse than
its selected checkpoint and worse than the base tokenizer-14 candidates.

\subsubsection{Compute Summary}

Table~\ref{tab:compute-summary} reports the logged wall-clock runtime recorded
in W\&B, together with GPU metadata from each run's \verb|wandb-metadata.json|
file. These values are approximate compute accounting numbers for the practical
xLSTM search: they combine runs made during different stages of the search and
should not be read as controlled hardware benchmarks.

\begin{table}[htbp]
	\centering
	\small
	\resizebox{\textwidth}{!}{%
		\begin{tabular}{lrrlll}
			\toprule
			Experiment stage                & Runs & Logged hours & GPU used                 & VRAM       & Purpose                                               \\
			\midrule
			Warm-up scale/batch checks      & 4    & 6.78         & RTX 3050 / RTX 3060      & 6 / 12 GiB & Verify training behavior and basic scale choices      \\
			Learning-rate sweep             & 10   & 15.86        & RTX 3060                 & 12 GiB     & Compare scheduler and learning-rate settings          \\
			Tokenizer sweep                 & 8    & 13.85        & RTX 3060                 & 12 GiB     & Compare REMI-style tokenizer conditions               \\
			Refined tokenizer--batch sweep  & 6    & 11.53        & RTX 3060                 & 12 GiB     & Recheck tokenizer 10/14 with batch sizes 4, 8, and 16 \\
			Medium-model follow-up          & 1    & 17.76        & RTX 3060                 & 12 GiB     & Test larger xLSTM capacity with longer training       \\
			\midrule
			Total summarized xLSTM training & 29   & 65.78        & Mixed RTX 3050/3060 runs & 6--12 GiB  & Validation-stage experimental budget                  \\
			\bottomrule
		\end{tabular}
	}
	\caption{Approximate compute used by completed xLSTM experiments, aggregated from W\&B runtime summaries and run metadata.}
	\label{tab:compute-summary}
\end{table}

The compute distribution explains why the early results should be treated as a
model-selection phase. Most of the budget was spent on validation sweeps rather
than on one final long training run. This was appropriate because the main
uncertainty was representation and optimization choice; once tokenizer 14 and
batch size 16 were selected, the final comparison could spend compute on a
matched xLSTM, GPT-2-style, and LLaMA-style evaluation.

\subsection{Matched Training and Test Results}

\subsubsection{Validation and Test Results}

The matched comparison changes the role of the results from xLSTM search to
architecture comparison. All three models use tokenizer 14, batch size 16, a
1024-token context, eight training epochs, and 19,432 optimizer updates. Under
these matched conditions, the xLSTM checkpoint gives the strongest held-out
next-token result. Its test loss is 1.3495 and its test perplexity is 3.8554,
compared with 1.8555 and 6.3948 for the LLaMA-style Transformer, and 1.9657 and
7.1399 for the GPT-2-style Transformer.

The validation trajectories in Figure~\ref{fig:matched-training-curves} and
Figure~\ref{fig:matched-validation-loss} also show different training behavior.
xLSTM reaches its best validation checkpoint at step 4000 and remains relatively
stable. GPT-2 improves more slowly and is selected at step 8000. LLaMA is the
stronger Transformer on held-out token metrics, but its curve overfits after its
step-4000 checkpoint. The next-token evidence therefore favors xLSTM for the
matched setup, with LLaMA as the stronger of the two Transformer baselines.

\subsubsection{Compute Summary}

The final comparison is more controlled than the earlier xLSTM search because
the models share the same update budget and data conditions. Table~\ref{tab:matched-compute-summary}
summarizes the compute metadata recorded in W\&B for the three matched runs.

\begin{table}[htbp]
	\centering
	\scriptsize
	\resizebox{\textwidth}{!}{%
		\begin{tabular}{llrrrrll}
			\toprule
			Model                   & W\&B run                                     & Logged runtime (h) & Updates & Updates/s & Batch & GPU              & Program                             \\
			\midrule
			xLSTM                   & \texttt{xlstm-tok14-batch16-sched-decay8000} & 3.83               & 19,432  & 1.41      & 16    & RTX 3060, 12 GiB & \path{src.models.xlstm.train}       \\
			GPT-2-style Transformer & \texttt{tok\_14\_batch\_16\_gpt2}            & 2.04               & 19,432  & 2.64      & 16    & RTX 3060, 12 GiB & \path{src.models.Transformer.train} \\
			LLaMA-style Transformer & \texttt{tok\_14\_batch\_16\_llama}           & 1.94               & 19,432  & 2.78      & 16    & RTX 3060, 12 GiB & \path{src.models.Transformer.train} \\
			\bottomrule
		\end{tabular}
	}
	\caption{W\&B compute summary for the matched training runs in the \texttt{duet-of-models} project. Runtime is reported from W\&B summary fields and GPU metadata from the W\&B run metadata.}
	\label{tab:matched-compute-summary}
\end{table}

The Transformer rows use the Hugging Face Trainer's recorded
\verb|train_runtime| values, 7348.80 seconds for GPT-2 and 6981.35 seconds for
LLaMA. The xLSTM trainer does not export an equivalent \verb|train_runtime|
summary field, so its row uses the W\&B run \verb|_runtime| value of 13783.34
seconds. The updates-per-second value for xLSTM is therefore an approximate
logged-run rate, while the Transformer values come directly from the Trainer
summary. All three runs used the same GPU class, the same batch size, the same
epoch budget, and the same 19,432 optimizer updates, so the table is useful for
practical compute accounting. It should not be read as a hardware-normalized
architecture speed benchmark, because the xLSTM and Transformer trainers use
different implementations and log runtime through different code paths.

\subsection{Generation Results}

\subsubsection{Generation Metrics}

The generated-MIDI evaluation gives a more mixed result than the next-token
test metrics. All three selected checkpoints generated 200 continuations, and
all generated outputs decoded and scored successfully. GPT-2 has the lowest FMD
score at 201.91, followed by xLSTM at 202.54 and LLaMA at 205.64. These values
are close enough that FMD alone should not be treated as a complete ranking.

The MusPy descriptor differences in
Table~\ref{tab:matched-generation-metrics} and
Figure~\ref{fig:matched-generation-metric-distributions} separate the models
more clearly. xLSTM has the best empty-beat-rate and polyphony-rate agreement,
which indicates stronger preservation of rhythmic density and simultaneous-note
structure. LLaMA gives the best pitch-entropy agreement, suggesting closer pitch
distribution statistics. GPT-2 has the best FMD and pitch-range difference, but
also the weakest pitch-entropy and empty-beat-rate differences. The generation
results therefore do not overturn the next-token result, but they show that
different metrics favor different musical properties.

\subsubsection{Compute Summary}

Generation time is the largest practical difference in the final evaluation.
The generated-MIDI file timestamps imply about 5.50 hours for GPT-2, 8.78 hours
for LLaMA, and 44.81 hours for xLSTM over the same 200-sample continuation set.
Because all three runs generated the same total continuation length, the
approximate throughputs are 191.9 tokens/s for GPT-2, 120.2 tokens/s for LLaMA,
and 23.5 tokens/s for xLSTM. The slower xLSTM decoding is consistent with the
token-by-token cost of recurrent xLSTM blocks, especially the sLSTM blocks, but
the thesis treats this as practical runtime evidence rather than a controlled
architecture benchmark.

